\ExamPart{Câu tự luận}{3}{}

\NQuestion Làm theo các yêu cầu sau:
\begin{SubQuestions}
\item Giải bất phương trình $\left(x^2-5x+6\right)\left(x^2-1\right)\le 0$.
\item Tìm tất cả các giá trị của \(m\) để phương trình $x^2-2(m+1)x+m^2+3=0$ có hai nghiệm phân biệt \(x_1,x_2\) thỏa mãn $x_1+x_2>2$.
\end{SubQuestions}

\NQuestion Trong mặt phẳng tọa độ \(Oxy\), cho ba điểm $A(1;2)$, $B(5;4)$, $C(3;-2)$.
\begin{SubQuestions}
\item Tính tọa độ các vectơ \(\overrightarrow{AB}\), \(\overrightarrow{AC}\).
\item Chứng minh rằng tam giác \(ABC\) vuông tại \(A\).
\item Tính diện tích tam giác \(ABC\).
\end{SubQuestions}

\NQuestion Trong mặt phẳng tọa độ \(Oxy\), cho đường tròn \((C)\) có tâm \(I(2;1)\) và đi qua điểm \(A(5;5)\).
\begin{SubQuestions}
\item Viết phương trình đường tròn \((C)\).
\item Viết phương trình tiếp tuyến của \((C)\) tại điểm \(A\).
\end{SubQuestions}

\NQuestion Trong mặt phẳng tọa độ \(Oxy\), cho tam giác \(ABC\) với $A(0;0)$, $B(6;0)$, $C(0;8)$.
Gọi \(M\) là trung điểm của cạnh \(BC\).
\begin{SubQuestions}
\item Tính tọa độ điểm \(M\).
\item Viết phương trình đường thẳng \(AM\).
\item Tính khoảng cách từ điểm \(B\) đến đường thẳng \(AM\).
\end{SubQuestions}